\section{杜教筛}\label{sec:du-sieve}

\subsection{作用与适用范围}

杜教筛对任意数论函数 $f$ 求前缀和 $S(n)=\sum_{i=1}^{n}f(i)$。适用于 $n$ 大到不能预处理 $1..n$ 的情形。递归只依赖 $O(\sqrt n)$ 个不同的 $\lfloor n/i\rfloor$，配合记忆化可做到亚线性。线性筛预处理前 $B\approx n^{2/3}$ 项后复杂度为 $O(n^{2/3})$，不做预处理为 $O(n^{3/4})$。

它不要求 $f$ 是积性函数；需要构造数论函数 $g$（$f*g$ 为狄利克雷卷积），使
\begin{itemize}[leftmargin=1.8em, itemsep=0.15em, topsep=0.2em]
    \item $\sum_{i=1}^{n}(f*g)(i)$ 能快速求值；
    \item $\sum_{i=1}^{n}g(i)$ 能快速求值，且 $g(1)$ 非零（模意义下可逆）。
\end{itemize}

\subsection{算法简述}

按
$$
g(1)S(n)=\sum_{i=1}^{n}(f*g)(i)-\sum_{i=2}^{n}g(i)S\left(\left\lfloor\frac{n}{i}\right\rfloor\right)
$$
递归，右侧用数论分块，$S$ 记忆化。通常取 $g(1)=1$；若 $g(1)\ne1$，最后除以 $g(1)$。常见构造是 $\mu*1=\epsilon$、$\varphi*1=\operatorname{id}$；例如 $f=\operatorname{id}^2\varphi$ 时取 $g=\operatorname{id}^2$，则 $f*g=\operatorname{id}^3$。

\subsection{伪代码}

记 $\pvar{pre}[n]$ 为线性筛预处理的 $S(n)$，$\pvar{mem}$ 为记忆化表：

\begin{algorithm}[H]
    \caption{杜教筛}
    \begin{algorithmic}[1]
        \Procedure{SumFG}{$n$}
            \State \Return $\sum_{i=1}^{n}(f*g)(i)$
        \EndProcedure
        \Procedure{SumG}{$n$}
            \State \Return $\sum_{i=1}^{n}g(i)$
        \EndProcedure
        \Procedure{Sum}{$n$}
            \If{$n \le B$}
                \State \Return $\pvar{pre}[n]$
            \EndIf
            \If{$n \in \pvar{mem}$}
                \State \Return $\pvar{mem}[n]$
            \EndIf
            \State $\pvar{ans} \gets \Call{SumFG}{n}$
            \State $l \gets 2$
            \While{$l \le n$}
                \State $r \gets n / (n / l)$
                \State $\pvar{ans} \gets \pvar{ans} - (\Call{SumG}{r}-\Call{SumG}{l-1})\cdot \Call{Sum}{n/l}$
                \State $l \gets r + 1$
            \EndWhile
            \State $\pvar{mem}[n] \gets \pvar{ans} / g(1)$
            \State \Return $\pvar{mem}[n]$
        \EndProcedure
    \end{algorithmic}
\end{algorithm}
